\section{The Probabilistic Audit}\label{sec:estimation}

The theory says that local validity composes. The audit turns that structural fact into an operational workflow for a deployed tree. Read through the iff of Section~\ref{sec:theorems}, the audit is the finite-sample estimator of the structural-approximation-error assumption: it measures how far the realized approximation error deviates from vanishing on the C1/C2/C3 subspace, which is the same thing as measuring how far the trained operator lies from the projection target. The audit samples a finite set of leaves, merge nodes, and re-summarized outputs, runs local oracle checks on those units, and transports the resulting estimates into a distortion or gap certificate.

For a realized document $x$, let $U(x)$ denote the finite population of \emph{all realized nodes} in its reduction tree: leaves, internal merges, and the root. The node-sampling law assigns each $u \in U(x)$ its own logged propensity $\pi_u \in (0,1]$. The root gets no special-casing in this node-population definition: it enters the sampled-node estimand through its own propensity, like every other realized node.

This node-population estimand is distinct from the always-available document-level supervision channel used in our simulations. We write
\[
y_{\mathrm{doc}} := f^\ast(x)
\]
for the full-document label, and
\[
y_{\mathrm{tree}}
\]
for the final-summary/tree-side oracle target or prediction. These coincide only under exact preservation or exact-control assumptions. Outside those regimes, the top/document loss and the sampled-node tree loss are different objects and should be reported separately.

\paragraph{The audit in four steps.}
\begin{enumerate}[leftmargin=1.5em]
    \item Sample units from the finite populations of leaves, merge nodes, and optional re-summarization nodes.
    \item Run the relevant local oracle checks for C1, C2, and C3 on those sampled units.
    \item Estimate two kinds of quantities: binary violation rates for diagnosis and mean distortion for certificate transport.
    \item Combine the distortion estimate with calibration, estimation, and clipping envelopes to obtain a finite-sample gap certificate.
\end{enumerate}

The audit reports two different objects. \emph{Violation rates} are binary diagnostics telling us which local law fails and how often. \emph{Distortion means} are numeric quantities that feed the gap theorems. Both are estimated from the same sampled nodes, but they play different roles. In the iff language of Section~\ref{sec:theorems}, violation rates answer ``where does the approximation error fail to be zero on the C1/C2/C3 subspace?'' and distortion means answer ``how far is the learned operator from the projection target in oracle units?'' Same sample, two readings; the bound in \S\ref{sec:discussion} uses both.

\begin{figure}[t]
\centering
\resizebox{0.97\linewidth}{!}{%
\begin{tikzpicture}[
    node distance=0.7cm and 0.7cm,
    box/.style={rectangle, draw, rounded corners=2pt, minimum width=2.5cm, minimum height=0.8cm, align=center, font=\footnotesize},
    arr/.style={-{Stealth[length=4pt]}, thick}
]
\node[box, fill=blue!8] (tree) {Raw document\\$\rightarrow$ compression tree};
\node[box, fill=green!8, right=of tree] (sample) {Sample leaves, merges,\\re-summaries};
\node[box, fill=yellow!10, right=of sample] (checks) {Run C1/C2/C3 checks\\and distortion queries};
\node[box, fill=orange!10, right=of checks] (ht) {Estimate violation rates\\and HT distortion mean};
\node[box, fill=red!8, right=of ht] (cert) {Transport through\\gap certificate};

\draw[arr] (tree) -- (sample);
\draw[arr] (sample) -- (checks);
\draw[arr] (checks) -- (ht);
\draw[arr] (ht) -- (cert);
\end{tikzpicture}%
}
\caption{End-to-end audit pipeline. The same sampled nodes support both interpretable local diagnostics (which law fails) and the distortion estimate used in the preference-gap certificate.}
\label{fig:audit-pipeline}
\end{figure}

\subsection{The Audit Protocol}

The displayed formulas below are the equal-probability versions. Under unequal inclusion probabilities, replace them with the Horvitz--Thompson (HT) weighted analogs introduced below.

\paragraph{Sampling Leaves (Sufficiency).}
We sample $n_\ell$ leaf blocks $\{b_i\}$ uniformly from the corpus. For each block, we generate its summary $g(b_i)$ and compare the oracle values:
\[
\hat{p}_{\text{suff}} := \frac{1}{n_\ell} \sum_{i=1}^{n_\ell} \One\left\{ \metric\big(\fstar(g(b_i)), \fstar(b_i)\big) > \tau \right\},
\]
where $\tau$ is a tolerance threshold (usually 0).

\paragraph{Sampling Triangle / Merge Consistency.}
Condition C3 requires a chain of equivalences. We sample $n_m$ merge edges and audit whichever link fits in memory. For a sampled merge node $u$ with children $u_L, u_R$, we can compare the disjoint path either to the raw concatenation,
\[
\hat{p}_{\text{assoc,raw}} := \frac{1}{n_m} \sum_{j=1}^{n_m} \One\left\{ \metric\!\big(\fstar(g(g(u_L^j) \concat g(u_R^j))),\; \fstar(u_L^j \concat u_R^j)\big) > \tau \right\},
\]
or to the joint-summary path,
\[
\hat{p}_{\text{assoc,joint}} := \frac{1}{n_m} \sum_{j=1}^{n_m} \One\left\{ \metric\!\big(\fstar(g(g(u_L^j) \concat g(u_R^j))),\; \fstar(g(u_L^j \concat u_R^j))\big) > \tau \right\}.
\]
At leaf boundaries, $u_L$ and $u_R$ are adjacent raw blocks whose concatenation often fits in context; at internal merges, they are stored summaries.

\paragraph{Sampling Stability (Idempotence).}
If the pipeline includes post-processing rounds, we sample $n_s$ summaries $z_k$ that have already passed through $g$ and check if re-summarization alters the oracle:
\[
\hat{p}_{\text{idem}} := \frac{1}{n_s} \sum_{k=1}^{n_s} \One\left\{ \metric\big(\fstar(g(z_k)), \fstar(z_k)\big) > \tau \right\}.
\]

Under equal-probability sampling within each audit stratum, these sample means are design-unbiased for the corresponding finite-population violation rates. They are primarily diagnostic: they tell us whether the problem is leaf retention, reprocessing drift, or merge failure.

\paragraph{Binary rates vs.\ distortion.}
The gap theorems in Section~\ref{sec:theorems} bound mean distortion, not binary violation rates. A run can have a small violation rate but still have large distortion on the failing units, or vice versa. The natural workflow is therefore to report both: the $\hat{p}$-rates for interpretability and a distortion mean for certificate transport.

\subsection{Scaling with Tree Size}

Let a realized hierarchy contain $N$ leaves and $M = N-1$ merges. A transparent union bound gives
\begin{equation}\label{eq:error_budget}
\Prob\big[\text{root violation}\big]
\le
\sum_{i=1}^{N} p_{\text{suff},i}
\;+\;
\sum_{j=1}^{M} p_{\text{assoc},j}
\;+\;
\sum_{r=1}^{R-1} p_{\text{idem},r},
\end{equation}
where each term denotes the true violation probability of the corresponding checked event and $R$ counts optional clean-up rounds. Under homogeneous upper bounds this reduces to $N p_{\text{suff}}^{\max}+M p_{\text{assoc}}^{\max}+(R-1)p_{\text{idem}}^{\max}$, so even tiny local failure probabilities accumulate over large trees.

\paragraph{From violation rates to $\Delta_R$.}
The bridge is the distortion estimator. With the ideal oracle, instantiating HT with $Y_i^\star=\metric(\fstar(s_i),\fstar(S(i)))$ gives an unbiased estimate of the audited unit-population mean $\mu_{\mathrm{dist}}^\star$. Operationally we usually use an accessible judge $f$, yielding $\hat{\mu}_{\mathrm{HT}}^{J}$ unbiased for $\mu_{\mathrm{dist}}^{J}$; Theorem~\ref{thm:e2e} then carries the difference between the ideal and accessible distortion means through the calibration envelope $B_{\mathrm{cal}}$. Concretely, $\hat{\mu}_{\mathrm{HT}}^{J}$ is the finite-sample estimator of the population distortion $\Delta_R$ from Section~\ref{sec:theorems}; the empirical quantity denoted $\hat\Delta_R$ in Section~\ref{sec:discussion} is this HT mean (plus the union-bound decomposition across leaf, merge, and idempotence audit units).

These statistics describe \emph{edge failure rates}, not correctness of an individual document. If a document has $N$ leaves and sufficiency failure rate $p_{\text{suff}}$, the probability that every leaf is correct is approximately $(1-p_{\text{suff}})^N \approx e^{-Np_{\text{suff}}}$. For large $N$, reliable root summaries therefore require small local error rates.

\subsection{Design-Based Estimation}

We cannot afford to query the oracle at every node, so we sample. This is a finite-population survey problem: the audit units form a population, we draw a random subset, and we want an unbiased estimate of the population mean. For document $x$, the relevant finite population is $U(x)$ itself, not just the subset that happened to be queried. The standard tool is the Horvitz--Thompson estimator, which reweights each sampled unit by the inverse of its inclusion probability. For queried unit $i$ with inclusion indicator $Z_i$, propensity $\pi_i$, and contribution $Y_i$:
\[
\hat{\mu}_{\mathrm{HT}} = \frac{1}{N}\sum_{i \in \mathcal{U}} \frac{Z_i}{\pi_i}Y_i.
\]
This estimator is conditionally unbiased for the population mean under the usual positivity and logged-propensity conditions.

\paragraph{Document labels vs.\ sampled-node estimation.}
The HT estimand above concerns the sampled-node population mean over $\mathcal{U}=\bigcup_x U(x)$. The full-document label $y_{\mathrm{doc}}$ is a separate complete-label channel and enters without IPW correction. In the shared-$g$ simulations below, we train the same readout on every realized node, but we evaluate and report two quantities separately: (i) the document-level top loss against $y_{\mathrm{doc}}$, and (ii) the sampled-node full-tree estimand, with naive, HT, and H\'ajek estimators computed from the logged node samples.

The practical bound stack decomposes target error into calibration, estimation, and clipping envelopes in objective units. Writing
\[
G^\ast := C_{\mathrm{meth}}\mu_{\mathrm{dist}}^{\star},\quad
G^J := C_{\mathrm{meth}}\mu_{\mathrm{dist}}^{J},\quad
G^E := C_{\mathrm{meth}}\hat{\mu}_{\mathrm{HT}}^{J},\quad
G^C := C_{\mathrm{meth}}\hat{\mu}_{\mathrm{dist}},
\]
we obtain
\[
|G^\ast-G^C| \le \underbrace{|G^\ast-G^J|}_{B_{\mathrm{cal}}} + \underbrace{|G^J-G^E|}_{B_{\mathrm{est}}} + \underbrace{|G^E-G^C|}_{B_{\mathrm{clip}}}.
\]
Plainly: $B_{\mathrm{cal}}$ is judge mismatch, $B_{\mathrm{est}}$ is sampling error, and $B_{\mathrm{clip}}$ is clipping-induced bias. The formal certificate keeps this exact four-piece shape: a local-law term $C_{\mathrm{meth}}\Delta_R$ plus calibration, estimation, and clipping addends. Its high-probability version uses the corresponding event budgets $\delta_{\mathrm{cal}}$, $\delta_{\mathrm{est}}$, and $\delta_{\mathrm{clip}}$ and combines them by a union bound.

\paragraph{Recommended reporting template.}
\begin{table}[htbp]
\centering\small
\begin{tabular}{p{0.17\linewidth}p{0.28\linewidth}p{0.44\linewidth}}
\toprule
Piece & Reported quantity & What it means \\
\midrule
Calibration envelope $B_{\mathrm{cal}}$ & Gold-vs-judge distortion mismatch on a held-out set & How much the accessible judge may differ from the target oracle on the audited population \\
Estimation envelope $B_{\mathrm{est}}$ & Finite-population radius for $|\mu_{\mathrm{dist}}^{J}-\hat{\mu}_{\mathrm{HT}}^{J}|$ & Sampling uncertainty from querying only a subset of nodes \\
Clipping envelope $B_{\mathrm{clip}}$ & Exact or conservative bias from clipped inverse-propensity weights & Additional bias introduced by variance-control choices in the estimator \\
Confidence composition & $\delta=\delta_{\mathrm{cal}}+\delta_{\mathrm{est}}+\delta_{\mathrm{clip}}$ & Final coverage statement for the reported certificate \\
\bottomrule
\end{tabular}
\caption{Minimal reporting template for non-vacuous certificates. The template shows where each source of slack comes from; a bare final bound hides that.}
\label{tab:reporting-template}
\end{table}

A robustness-to-adversarial-sampling analysis (a constrained-design variance bound that shows why $\pi_{\min}$ is the primary robustness lever and why certificates should be reported over a stress-test grid of $(\pi_{\min}, w_{\max})$ values) is in Appendix~\ref{app:audit-robustness}. The same layer also gives deterministic clipped-vs-unclipped H\'ajek envelopes, so $B_{\mathrm{clip}}$ can be reported as an explicit sensitivity term rather than folded into the sampling radius.

\subsection{From Audit to Optimization}

Because the audit identifies \emph{which} edges fail and \emph{which} condition each failure violates, it also provides supervision for prompt engineering. In practice, one can treat leaf and merge prompts as parameters $\theta$ and update them against the violation signal. The same audit thus supports both certification and tuning.

\paragraph{Empirical supervision split.}
In the full-tree shared-$g$ experiments, every document supplies a complete document-level label at the top, while only a random sampled subset of realized nodes in $U(x)$ receives node-level feedback. In that node law the root has its own logged propensity when it is part of the sampled-node channel, like any other node. Held-out evaluation therefore reports the exact full-node objective, the sampled-node naive/HT/H\'ajek estimates, and the separate document-vs-final-summary top loss.

Adaptive chunking and learned judges introduce interference paths. To keep evaluation valid, reporting must be honest about adaptation roles (chunker, summarizer, oracle): adaptation on train-role partitions, evaluation on held-out role intersections. Full details of the honesty protocol are deferred to the companion paper.
